\chapter{Torsional Memory and Bounce Conditions in Einstein--Cartan}
\label{ch:rebote-ec}

\section{Transport of the Perron law to the spin spectral density}

The realization in Perron cosmological form constructs two sequences on the
same closed route:

\begin{equation}
 \frac{a_n}{a_0}=\varphi^{n/3},
 \qquad
 M_n=\varphi^{-2n}.
 \label{eq:rebote-perron-pair}
\end{equation}

Elimination of level \(n\) produces the exact identity

\begin{equation}
 \boxed{M_n=\left(\frac{a_n}{a_0}\right)^{-6}.}
 \label{eq:rebote-memory-six}
\end{equation}

The exponent has a definite genealogy: the spatial dimension
contributes the factor \(3\), and the quadratic character of stable memory
contributes the factor \(2\).  The power \(6=2\cdot3\) thus expresses the dilution
of a quadratic density associated with an extensive spin charge.

\begin{definition}[Torsional transducer]
Having fixed a positive amplitude \(s_*^2\) on an initial section, the torsional
transducer is
\begin{equation}
 \mathcal T_{\rm EC}:M_n\longmapsto
 s_n^2=s_*^2M_n
 =s_*^2\left(\frac{a_n}{a_0}\right)^{-6}.
 \label{eq:rebote-transductor}
\end{equation}
Its domain is the stable Perron mode, and its codomain is the
quadratic spin sector of an Einstein--Cartan realization.
\end{definition}

The equation \cref{eq:rebote-transductor} preserves the exponent and
positivity.  The amplitude \(s_*^2\), the sign with which the contorsion enters
the gravitational equation, and the dimensional normalization belong to the
physical realizer.  This separation establishes exactly what the
HMT trunk supplies and what Cartan dynamics must supply.

\section{Homogeneous Equations and Density Convention}

To present the structural mechanism, we temporarily adopt units
\(c=1\) and a spatially flat FRW geometry.  Let

\begin{equation}
 \rho_m(a)=\rho_0a^{-3(1+w)},
 \qquad
 \rho_s(a)=s_0a^{-6},
 \label{eq:rebote-densities}
\end{equation}

with \(\rho_0>0\), \(s_0>0\), and \(w\) constant.  The contribution from
contorsion enters with a repulsive sign in the effective equation:

\begin{equation}
 H^2=\frac{8\pi G}{3}\bigl(\rho_m-\rho_s\bigr).
 \label{eq:rebote-friedmann}
\end{equation}

If energy densities are used, the right-hand side is multiplied by
\(c^{-2}\); if mass densities are used, it retains the form
\cref{eq:rebote-friedmann}.  The document maintains a single convention in
each calculation and restores the dimensional chart at the end.

The corresponding Raychaudhuri equation is

\begin{equation}
 \dot H=-4\pi G\bigl[(1+w)\rho_m-2\rho_s\bigr].
 \label{eq:rebote-raychaudhuri}
\end{equation}

The second contribution behaves as a stiff fluid with equation of
state \(w_s=1\) and negative effective density. This form is the homogeneous
publication of torsional repulsion at high density.

\section{Existence and uniqueness of the bounce radius}

\begin{theorem}[Regular reflection for \(w<1\)]
Assume \(\rho_0>0\), \(s_0>0\), and \(w<1\).  Then there exists a unique
positive radius \(a_b\) such that \(H(a_b)=0\), given by
\begin{equation}
 \boxed{
 a_b^{3(1-w)}=\frac{s_0}{\rho_0}.}
 \label{eq:rebote-radius}
\end{equation}
At that point
\begin{equation}
 \boxed{
 \dot H_b=4\pi G(1-w)\rho_b>0,}
 \qquad
 \rho_b:=\rho_m(a_b)=\rho_s(a_b),
 \label{eq:rebote-positive-derivative}
\end{equation}
so that the contractive branch is extended into a regular expansive branch.
\end{theorem}

\begin{proof}
The condition \(H=0\) is equivalent to
\(\rho_0a^{-3(1+w)}=s_0a^{-6}\).  Multiplication by \(a^6\) yields
\cref{eq:rebote-radius}.  The exponent \(3(1-w)\) is positive; therefore the
power defines a bijection of \(\RR_{>0}\) onto itself and the radius is
unique.  Substituting \(\rho_m=\rho_s=\rho_b\) into
\cref{eq:rebote-raychaudhuri},
\[
 \dot H_b=-4\pi G[(1+w)-2]\rho_b
 =4\pi G(1-w)\rho_b>0.
\]
Positivity turns the simple zero of \(H\) into a minimum of \(a(t)\).
\end{proof}

The local expansion around \(t_b\) makes the regularity explicit:

\begin{equation}
 a(t)=a_b\left[
 1+2\pi G(1-w)\rho_b(t-t_b)^2
 +O\!\left((t-t_b)^3\right)
 \right].
 \label{eq:rebote-local-expansion}
\end{equation}

The law \(a^{-6}\) supplies the power that dominates the material sector when
\(w<1\) and \(a\) decreases. The torsional sign and positive amplitude cause
that dominance to halt the contraction before reaching \(a=0\).

\section{Limit Cases of Torsional Dynamics and Obstructions to the Regularity of the Bounce}

The value \(w=1\) produces the same exponent \(a^{-6}\) in both terms. The difference then retains a fixed sign determined by \(\rho_0-s_0\); pointwise equality ceases to select an isolated radius. For \(w>1\), the material density grows faster than the torsional term during contraction, and the preceding mechanism does not guarantee a bounce. These cases provide internal negative controls: the closure depends on the inequality of the exponents and not on a retrospective choice of the threshold.

Spatial curvature \(k\), a cosmological constant, and Bianchi anisotropies
add terms to \cref{eq:rebote-friedmann}. Their local effect
is controlled without presupposing the position of the new zero. Define
\[
 F_0(a):=\frac{8\pi G}{3}\bigl(\rho_m(a)-\rho_s(a)\bigr).
\]
The perturbed equation takes the form
\(H^2=F_R(a):=F_0(a)+R(a)\).

\begin{theorem}[Persistence \(C^1\) of the rebote transversal]
\label{thm:rebote-c1-persistence}
Let \(a_b\) be the radius of \cref{eq:rebote-radius}. There exist
\(0<a_-<a_b<a_+\) such that, for
\(I=[a_-,a_+]\),
\[
 F_0(a_-)<0<F_0(a_+),
 \qquad
 m_I:=\inf_{a\in I}F_0'(a)>0.
\]
If \(R\in C^1(I)\) satisfies
\begin{equation}
 \|R\|_{C^0(I)}
 <\delta_I
 :=\min\{-F_0(a_-),F_0(a_+)\},
 \qquad
 \|R'\|_{C^0(I)}<m_I,
 \label{eq:rebote-c1-bounds}
\end{equation}
then there exists a unique \(a_R\in(a_-,a_+)\) with \(F_R(a_R)=0\), and
that zero is transversal:
\begin{equation}
 F_R'(a_R)>0,
 \qquad
 \dot H_R(a_R)=\frac{a_R}{2}F_R'(a_R)>0.
 \label{eq:rebote-perturbed-transversality}
\end{equation}
For a family \(R(a,\eta)\) of class \(C^1\), the implicit function
theorem provides a local branch \(a_R(\eta)\) of class \(C^1\).
\end{theorem}

\begin{proof}
The identity \cref{eq:rebote-positive-derivative} is equivalent to \(F_0'(a_b)>0\). By continuity, \(I\) can be chosen with \(m_I>0\) and with the indicated change of sign. The first bound for \cref{eq:rebote-c1-bounds} preserves \(F_R(a_-)<0<F_R(a_+)\); the intermediate value theorem produces a zero. The second bound gives \(F_R'(a)\geq m_I-\|R'\|_{C^0(I)}>0\), so \(F_R\) is strictly increasing and the zero is unique. On each branch with \(H\neq0\), differentiating \(H^2=F_R(a)\) and using \(\dot a=aH\) gives \(\dot H=(a/2)F_R'(a)\); continuity extends the equality to \(a_R\). The nondegeneracy \(F_R'(a_R)>0\) is precisely the hypothesis of the implicit function theorem.
\end{proof}

\section{Torsional amplitude from the Cartan screen}

The memory exponent is closed by
\cref{eq:rebote-memory-six}. The amplitude may be sought prospectively in
the screen response. Let the exact identity be

\begin{equation}
 (\beta_m^{\square})^*\Lambda_m^{\square}\beta_m^{\square}
 =\frac{56}{81}I_{Q_{\square}}.
 \label{eq:rebote-screen-energy}
\end{equation}

A complete gravitational realization must construct a spin current

\begin{equation}
 \mathcal J_m
 =\mathcal J(\beta_m^{\square},U_m,c_s,c_g)
 \label{eq:rebote-spin-current}
\end{equation}

and a positive functional

\begin{equation}
 s_{0,m}=\mathfrak S(\mathcal J_m,\Lambda_m^{\square})
 \label{eq:rebote-amplitude-functional}
\end{equation}

compatible with refinement.  The pertinent nonadic normalization is the
tracial expectation, because it preserves the density of
\cref{eq:rebote-screen-energy}; the raw sum counts nine copies and belongs
to a distinct extensive quantity.

\begin{proposition}[Obligations of the amplitude selector]
An admissible selector for \(s_0\) must satisfacer:
\begin{enumerate}[label=\roman*)]
 \item positivity and dimensions of squared spin density;
 \item gauge invariance under the reduction induced by the transported cubic structure;
 \item compatibility with edge inversion and carry;
 \item naturality under \(9\)-ary refinement;
 \item independence from any bounce radius used as a target value.
\end{enumerate}
\end{proposition}

These conditions turn the amplitude into a localized algebraic constraint. The screen-energy equality supplies the quadratic form; the Cartan realizer must determine how that form contracts with the physical current.

\section{Torsion, orientation, and transport memory}

The metric publishes distances and causal cones.  The Cartan connection
also publishes the orientation of transport.  In HMT, route, carry, and
deep state already constitute genealogical data.  The correspondence

\begin{equation}
 \text{stable route memory}
 \longrightarrow
 \text{spin current}
 \longrightarrow
 \text{contorsion}
 \longrightarrow
 \rho_s\propto a^{-6}
 \label{eq:rebote-genealogy}
\end{equation}

provides a physical reading of that oriented information. Curvature
describes the response of geometry to energy; torsion preserves
the history of orientation that accompanies transport. The sixth-order
law is the point at which genealogical memory and quadratic spin
density have exactly the same scaling covariance.

\section{Mach's principle and cosmographic autoinertia}

The local state is evaluated within a route compatible with the global
boundary. The inertial response therefore belongs to the holonomic
stabilizer of the totality. In the homogeneous sector, the scale \(a_n\), the
memory \(M_n\), and the torsional current are three readings of the same enriched
genealogical state:

\begin{equation}
 X_n\longmapsto
 \bigl(a_n,M_n,\mathcal J_n\bigr).
 \label{eq:rebote-mach-triple}
\end{equation}

The Machian interpretation thereby acquires operatorial content: the local response depends on the global closure class, and the observer forms part of the prolonged section. Self-inertia is the preservation of the relative state under the holonomy that jointly transports geometry, observer, and memory.

\section{Exact identities of torsional dynamics and Einstein–Cartan interface conditions}

\begin{proofstatusbox}
The law \(M_n=(a_n/a_0)^{-6}\) is an exact internal theorem. The bounce theorem \cref{eq:rebote-radius,eq:rebote-positive-derivative} is an exact consequence of the homogeneous Einstein--Cartan equations \cite{Hehl1976} with \(\rho_0,s_0>0\) and \(w<1\). The transducer \cref{eq:rebote-transductor} exactly preserves the scaling law. The determination of \(s_0\) from the screen current, the extension to anisotropies, and the complete derivation of the Einstein--Cartan action are delimited physical interfaces.
\end{proofstatusbox}

\begin{falsifierbox}
The internal block is refuted if the Perron mode violates \cref{eq:rebote-memory-six}. The Einstein--Cartan realization is refuted if the induced current produces another power, a nonpositive amplitude, or a sign that eliminates \cref{eq:rebote-positive-derivative}. The homogeneous bounce fails for the declared data if the radius \cref{eq:rebote-radius} does not exist, if \(H(a_b)\ne0\), or if \(\dot H_b\le0\). The complete promotion is rejected if the amplitude functional retrospectively consults the desired radius or loses gauge, carry, or refinement.
\end{falsifierbox}
